\documentclass[10pt,a4paper]{amsart}
\usepackage[margin=19mm]{geometry}
\usepackage{amsmath,amssymb,mathtools}
\usepackage[scr=rsfs,cal=euler]{mathalfa}
\usepackage{xfrac}
\usepackage[unicode,hidelinks,bookmarks=false]{hyperref}
\newcommand{\ie}{i.e.\ }
\newcommand{\cf}{cf.\ }
\newcommand{\resp}{respectively\ }
\newcommand{\Subsection}{\subsection}
\providecommand{\nobreakdash}{\nobreak\hbox{-}}
\setlength{\emergencystretch}{2em}
\multlinegap0pt
\usepackage{amssymb}
\usepackage{mathtools}
\usepackage{xcolor}
\usepackage[shortlabels]{enumitem}
\usepackage{tikz}
\usepackage{float}
\newtheorem{lemma}{Lemma}
\newtheorem{proposition}{Proposition}
\def\l{\langle}
\def\r{\rangle}
\title{SIMIS and packing properties of Alexander dual of connected ideals}
\author{Om Prakash Bhardwaj, Kanoy Kumar Das,, Rutuja Sawant}
\date{Source: arXiv:2603.18751v1 (CC BY 4.0)}
\begin{document}
\maketitle

\noindent\emph{Source and licence.} SIMIS and packing properties of Alexander dual of connected ideals. Version arXiv:2603.18751v1 (CC BY 4.0), distributed by arXiv under the Creative Commons Attribution 4.0 International licence, http://creativecommons.org/licenses/by/4.0/, as recorded in the arXiv licence metadata for this exact version and shown on the abstract page https://arxiv.org/abs/2603.18751v1. Full licence notice: \texttt{licenses/arxiv-2603.18751-CC-BY-4.0.txt}.\par\smallskip

\noindent\textit{Editorial scope.} Counted unit: the article's proposition prop:symbolicpowers-cycles (source0.tex lines 455--469). The article's complete original proof of it is delivered with it (source0.tex lines 470--592, one \texttt{proof} environment of 123 lines). No mathematics is written, repaired or re-derived: the statement and the proof are the article's own lines, in the article's own notation, and no retained line is substituted or re-flowed.  Retained uncounted as the source-local context the proof needs: lines 213--215; lines 259--268; lines 275--277; lines 325--327.  Nothing was omitted from any retained span, so the package carries no omission record.  No retained line contains a citation, so the wrapper carries no bibliography and no bibliography entry.  No retained line contains a figure environment, an image-inclusion command or drawing code, so no image is delivered and the package declares no asset dependency.  Editorial formatting only, in addition: an AMS \texttt{amsart} preamble that restates the article's own declarations and macros used by the retained lines, namely the environments \texttt{lemma}, \texttt{proposition} on separate counters, together with the standard packages those lines need.  The retained environments print in this document as Lemma 1, Proposition 1 in order of appearance, whereas the article prints the same items with the numbering of its own full document.  The complete native source of the article as distributed in the arXiv e-print for this exact version is bundled unchanged as \texttt{sources/196716/source0.tex}.
\begin{lemma}\label{lem: Simis implies packed}
    Let $I\subseteq R$ be a square-free monomial ideal. If $I$ is Simis then $I$ satisfies the packing property. 
\end{lemma}

\begin{proposition}\label{prop: mingens of J_t(P_n)}
    For $1\leq i_1<i_2<\cdots < i_r\leq n$, $\mathbf{x}^{\alpha}=x_{i_1}\cdots x_{i_r}\in \mathcal{G}(J_t(P_n))$ if and only if the following holds:
    \begin{enumerate}
        \item $r\geq \lfloor \frac{n}{t} \rfloor$;
        \item $i_1\leq t, i_2\geq t+1$;
        \item $i_r\geq n-t+1, i_{r-1}\leq n-t$;
        \item $i_{j+1}-i_{j}\leq t$ for all $1\leq j \leq r-1$;
        \item $i_{j+2}-i_{j}\geq t+1$ for all $1\leq j\leq r-2$.
    \end{enumerate}
\end{proposition}

\begin{proposition}\label{prop: Symbolic Prop for P_n}
    For $n\geq t\geq 2$, we have $J_t(P_n)^{(s)}= J_t(P_n)^{s}$ for all $s\geq 1$. In particular, $J_t(P_n)$ has the packing property. 
\end{proposition}

\begin{lemma}\label{lem:cycle-not-packed}
    Let $n>t$ be positive integers such that $(n,t) \notin \{(6,3),(9,3), (8,4)\}$. Then $J_t(C_n)$ is not packed.    
\end{lemma}

\begin{proposition}\label{prop:symbolicpowers-cycles}
    The following statements are true:
    \begin{enumerate}
        \item $J_t(C_n)^{(s)}= J_t(C_n)^{s}$ for all $s\geq 1$ if and only if 
\begin{enumerate}
            \item $t=3, n=3,6,9$;
            
            \item $t=4, n=4,8$;
            
            \item for $t\geq 5$, $n=t$.
        \end{enumerate}
        
        \item $J_t(C_n)^{(s)}= J_t(C_n)^{s}$ for all $s\geq 1$ if and only if $J_t(C_n)$ is packed.
    \end{enumerate}    
\end{proposition}

\begin{proof}
    \medskip \noindent
    (1) For the `if' part, observe that when $n=t$, the ideal $J_t(C_t)$ is generated by variables. Consequently, $J_t(C_t)^{(s)} = J_t(C_t)^s$ for all $s \ge 1$. Now assume that $n \ne t$. We treat the remaining cases for $n$ and $t$ separately.

    
    We first consider the case $n=6$ and $t=3$. We show that $J_3(C_6)^{(s)} = J_3(C_6)^s$ for all $s \ge 1$. The proof proceeds by induction on $s$. The assertion is clear for $s=1$. Assume that $J_3(C_6)^{(\ell)} = J_3(C_6)^{\ell}$ for every $\ell < s$. It therefore suffices to prove that $J_3(C_6)^{(s)} \subseteq J_3(C_6)^s.$ Let $f \in J_3(C_6)^{(s)}$. Without loss of generality, assume that $x_6 \mid f$. From the description of symbolic powers, it follows that
    \[
    J_3(C_6)^{(s)} =\l x_4,x_5,x_6 \r^s
    \cap \l x_5,x_6,x_1\r^s
    \cap \l x_6,x_1,x_2 \r^s
    \cap J_3(P_5)^{(s)}.
    \]
    In particular, $f \in J_3(P_5)^{(s)}$. Hence, by Proposition~\ref{prop: Symbolic Prop for P_n}, there exists $g \in \mathcal{G}(J_3(P_5))$ such that $g \mid f$ and $\frac{f}{g} \in J_3(P_5)^{s-1}= J_3(P_5)^{(s-1)}.$ We consider all possible choices of the generator $g$.
    \begin{enumerate}
        \item If $g = x_1x_4$ or $g = x_2x_5$, then $g \in \mathcal{G}(J_3(C_6))$ and $g \mid f$. Note that $\frac{f}{g} \in J_3(C_6)^{(s-1)} = J_3(C_6)^{s-1}$, where the last equality follows from the induction hypothesis. Hence $f \in J_3(C_6)^s$.

        \item If $g = x_3$, set $g' = x_3x_6$. Then $g' \in \mathcal{G}(J_3(C_6))$ and $g' \mid f$. Again, observe that $\frac{f}{g'} \in J_3(C_6)^{(s-1)} = J_3(C_6)^{s-1}$, where the last equality comes from the induction hypothesis. Thus $f \in J_3(C_6)^s$.

        \item If $g = x_2x_4$, we may assume that $x_1 \nmid f$, $x_3 \nmid f$, and $x_5 \nmid f$; otherwise one of the previous cases would apply. Hence $x_6^s \mid f$. If we set $g' = x_2x_4x_6$, then observe that $\frac{f}{g'} \in J_3(C_6)^{(s-1)} = J_3(C_6)^{s-1}$, and therefore $f \in J_3(C_6)^s$.
    \end{enumerate}
    Thus $J_3(C_6)^{(s)} \subseteq J_3(C_6)^s$, and hence, $J_3(C_6)^{(s)} = J_3(C_6)^s$ for all $s \ge 1$.


Next, we consider the case $n=9$ and $t=3$. To show $J_3(C_9)^{(s)} = J_3(C_9)^s$ for all $s \ge 1$, we again proceed by induction on $s$. The statement is true for $s=1$. Assume that $J_3(C_9)^{(\ell)} = J_3(C_9)^{\ell}$ for every $\ell<s$. It now suffices to prove that $J_3(C_9)^{(s)} \subseteq J_3(C_9)^s.$ Let $f\in J_3(C_9)^{(s)}$ and assume $x_9\mid f$. From the expression of symbolic powers, we get
\[
J_3(C_9)^{(s)} = \l x_7,x_8,x_9 \r ^s
\cap \l x_8,x_9,x_1 \r ^s
\cap \l x_9,x_1,x_2\r ^s
\cap J_3(P_8)^{(s)}.
\]
Hence, by Proposition~\ref{prop: Symbolic Prop for P_n}, there exists
$g\in \mathcal{G}(J_3(P_8))$ such that
$g\mid f$ and $\frac{f}{g}\in J_3(P_8)^{s-1}
=J_3(P_8)^{(s-1)}.$ By Proposition~\ref{prop: mingens of J_t(P_n)}, we have
\begin{align*}   
\mathcal{G}(J_3(P_8)) = \left\lbrace
\begin{array}{c}
    x_1x_4x_5x_8, x_1x_4x_6, x_1x_4x_7, x_2x_4x_6,x_2x_4x_7,x_2x_5x_6, \\ [1mm]
    x_2x_5x_7,x_2x_5x_8,x_3x_4x_7,x_3x_5x_7,x_3x_5x_8,x_3x_6
\end{array}
\right\rbrace .
\end{align*}
We now consider the possible choices for the generator $g$ case by case.

\begin{enumerate}

\item If $g=x_1x_4x_7$ or $g=x_2x_5x_8$, then
$g\in \mathcal{G}(J_3(C_9))$.
Thus $g\mid f$ and moreover,
$\frac{f}{g}\in J_3(C_9)^{(s-1)}=J_3(C_9)^{s-1}$, where the last equality is by induction. Therefore, $f\in J_3(C_9)^s$.

\item If $g=x_3x_6$, set $g'=x_3x_6x_9$.
Then $g'\in \mathcal{G}(J_3(C_9))$ and $g'\mid f$.
Note that $\frac{f}{g'}\in J_3(C_9)^{(s-1)}$, and hence by the induction hypothesis, $\frac{f}{g'}\in J_3(C_9)^{s-1}$. Therefore, $f\in J_3(C_9)^s$.

\item If $g=x_1x_4x_6$, we may assume
$x_3\nmid f$ and $x_7\nmid f$; otherwise one of the previous cases applies. It then follows that
$x_1^{\alpha_1}x_2^{\alpha_2}\mid f$ with $\alpha_1+\alpha_2\ge s$ and $x_8^{\beta_1}x_2^{\beta_2}\mid f$ with $\beta_1+\beta_2\ge s$. Let $g'=x_1x_4x_6x_9\in J_3(C_9)$. Then observe that $\frac{f}{g'}\in J_3(C_9)^{(s-1)}=\frac{f}{g'}\in J_3(C_9)^{s-1}$, where the equality follows by induction. Hence, $f\in J_3(C_9)^s$. A similar argument applies to the cases when  $g\in \{ x_3x_5x_8, x_2x_4x_6, x_3x_5x_7, x_2x_5x_6, x_3x_4x_7, x_2x_4x_7, x_2x_5x_7\}$.

\item If $g=x_1x_4x_5x_8$, we may assume that $x_2,x_3, x_6,x_7\nmid f$; otherwise, one of the previous cases applies. It then follows that $x_1^s\mid f$ and $x_8^s\mid f$. Hence, $\frac{f}{g}\in J_3(C_9)^{(s-1)}=J_3(C_9)^{s-1}$, where the equality follows from the induction hypothesis.
Since $g\in \mathcal{G}(J_3(C_9))$, we conclude that
$f\in J_3(C_9)^s$.

\end{enumerate}
Thus, in any case $f\in J_3(C_9)^s$, and therefore $J_3(C_9)^{(s)} = J_3(C_9)^s$ for all $s\ge1$.

Finally, consider the case $n=8$ and $t=4$. We prove that
$J_4(C_8)^{(s)} = J_4(C_8)^s$ for all $s \ge 1$ by induction on $s$. Assume that $J_4(C_8)^{(\ell)} = J_4(C_8)^{\ell}$ for every $1\le \ell < s$. Without loss of generality, assume that $f \in J_4(C_8)^{(s)}$ with $x_8 \mid f$. From the expressions of symbolic powers, we have
\[
J_4(C_8)^{(s)}= \l x_5,x_6,x_7,x_8\r ^s
\cap \l x_6,x_7,x_8,x_1\r ^s
\cap \l x_7,x_8,x_1,x_2\r ^s
\cap \l x_8,x_1,x_2,x_3\r ^s
\cap J_4(P_7)^{(s)}.
\]
Since $f \in J_4(P_7)^{(s)}$, by Proposition~\ref{prop: Symbolic Prop for P_n}, $f \in J_4(P_7)^{s}$, and hence, there exists $g \in \mathcal{G}(J_4(P_7))$ such that $g \mid f$ and $\frac{f}{g} \in J_4(P_7)^{s-1}$. By Proposition~\ref{prop: mingens of J_t(P_n)}, we have 
\[
\mathcal{G}(J_4(P_7))=\{x_1x_5,\; x_2x_5,\; x_2x_6,\; x_3x_5,\; x_3x_6,\; x_3x_7,\; x_4\}.
\]
We now consider the possible choices for the generator $g$ case by case.
\begin{enumerate}


\item If $g \in \{x_1x_5, x_2x_6, x_3x_7\}$, then
$g \in \mathcal{G}(J_4(C_8))$. Observe that $\frac{f}{g} \in J_4(C_8)^{(s-1)} = J_4(C_8)^{s-1}$, where the equality follows from the induction hypothesis. Hence $f \in J_4(C_8)^s$.

\item If $g = x_4$, take $g' = x_4x_8$.
Then $g' \in \mathcal{G}(J_4(C_8))$ and $g' \mid f$.
Moreover, $\frac{f}{g'} \in J_4(C_8)^{(s-1)} = J_4(C_8)^{s-1}$,
so $f \in J_4(C_8)^s$.

\item If $g = x_2x_5$, we may assume that
$x_1 \nmid f$, $x_4 \nmid f$, and $x_6 \nmid f$;
otherwise one of the previous cases applies.
Then $x_5^{\alpha_1}x_7^{\alpha_2} \mid f$ with $\alpha_1+\alpha_2 \ge s$,
$x_7^{\beta_1}x_8^{\beta_2} \mid f$ with $\beta_1+\beta_2 \ge s$,
and $x_2^{\gamma_1}x_3^{\gamma_2} \mid f$ with $\gamma_1+\gamma_2 \ge s$.
Let $g' = x_2x_5x_8$. Then
$\frac{f}{g'} \in J_4(C_8)^{(s-1)} = J_4(C_8)^{s-1}$,
and hence $f \in J_4(C_8)^s$.

\item If $g = x_3x_5$, we may assume that
$x_1,x_2,x_4,x_7 \nmid f$; otherwise a previous case applies.
Then $x_8^s \mid f$.
Let $g' = x_3x_5x_8$. Hence
$\frac{f}{g'} \in J_4(C_8)^{(s-1)} = J_4(C_8)^{s-1}$,
and therefore $f \in J_4(C_8)^s$.

\item If $g = x_3x_6$, we may assume that
$x_2,x_4,x_5,x_7 \nmid f$; otherwise one of the earlier cases applies.
Thus $x_3^s \mid f$ and $x_6^s \mid f$.
Let $g' = x_3x_5x_8$. Then
$\frac{f}{g'} \in J_4(C_8)^{(s-1)} = J_4(C_8)^{s-1}$,
so $f \in J_4(C_8)^s$.

\end{enumerate}
Therefore, in any case $f\in J_4(C_8)^s$. Hence, we can conclude that $J_4(C_8)^{(s)} = J_4(C_8)^s$ for all $s \ge 1$. This completes the proof for the `if' part. The `only if' part follows from Lemma~\ref{lem:cycle-not-packed}
and Lemma~\ref{lem: Simis implies packed}.

\medskip 
\noindent
(2) Follows from (1) and Lemma~\ref{lem:cycle-not-packed}.
\end{proof}

\end{document}
